\documentclass[10pt,a4paper]{amsart}
\usepackage[margin=19mm]{geometry}
\usepackage{amsmath,amssymb,mathtools}
\usepackage[scr=rsfs,cal=euler]{mathalfa}
\usepackage{xfrac}
\usepackage[unicode,hidelinks,bookmarks=false]{hyperref}
\newcommand{\ie}{i.e.\ }
\newcommand{\cf}{cf.\ }
\newcommand{\resp}{respectively\ }
\newcommand{\Subsection}{\subsection}
\providecommand{\nobreakdash}{\nobreak\hbox{-}}
\setlength{\emergencystretch}{2em}
\multlinegap0pt
\usepackage{mathtools}
\usepackage{amssymb}
\usepackage{dsfont}
\newtheorem{thm}{Theorem}
\renewcommand{\sinh}{\operatorname{sinh}}
\title{Dirichlet problems and exit distributions for the telegraph process and its planar extensions}
\author{Manfred Marvin Marchione, Enzo Orsingher}
\date{Source: arXiv:2605.05430v1 (CC BY 4.0)}
\begin{document}
\maketitle

\noindent\emph{Source and licence.} Dirichlet problems and exit distributions for the telegraph process and its planar extensions. Version arXiv:2605.05430v1 (CC BY 4.0), distributed by arXiv under the Creative Commons Attribution 4.0 International licence, http://creativecommons.org/licenses/by/4.0/, as recorded in the arXiv licence metadata for this exact version and shown on the abstract page https://arxiv.org/abs/2605.05430v1. Full licence notice: \texttt{licenses/arxiv-2605.05430-CC-BY-4.0.txt}.\par\smallskip

\noindent\textit{Editorial scope.} Counted unit: the article's thm tildeuthm (source0.tex lines 523--529). The article's complete original proof of it is delivered with it (source0.tex lines 530--551, one \texttt{proof} environment of 22 lines). No mathematics is written, repaired or re-derived: the statement and the proof are the article's own lines, in the article's own notation, and no retained line is substituted or re-flowed.  Retained uncounted as the source-local context the proof needs: lines 481--486; lines 509--514; lines 519--519.  Nothing was omitted from any retained span, so the package carries no omission record.  No retained line contains a citation, so the wrapper carries no bibliography and no bibliography entry.  No retained line contains a figure environment, an image-inclusion command or drawing code, so no image is delivered and the package declares no asset dependency.  Editorial formatting only, in addition: an AMS \texttt{amsart} preamble that restates the article's own declarations and macros used by the retained lines, namely the environments \texttt{thm} on separate counters, together with the standard packages those lines need.  The retained environments print in this document as Theorem 1 in order of appearance, whereas the article prints the same items with the numbering of its own full document.  The complete native source of the article as distributed in the arXiv e-print for this exact version is bundled unchanged as \texttt{sources/199950/source0.tex}.
\begin{equation}\label{algebrain}
\begin{aligned}
\widetilde{u}_0(\alpha,y\mathord{;}z)=&\frac{\lambda}{2(\lambda+ic\alpha)}\;\big[\widetilde{u}_1(\alpha,y\mathord{;}z)+\widetilde{u}_3(\alpha,y\mathord{;}z)\big]\\
\widetilde{u}_2(\alpha,y\mathord{;}z)=&\frac{\lambda}{2(\lambda-ic\alpha)}\;\big[\widetilde{u}_1(\alpha,y\mathord{;}z)+\widetilde{u}_3(\alpha,y\mathord{;}z)\big].
\end{aligned}
\end{equation}

\begin{equation}\label{bivFouriersys2}
\begin{cases}
\dfrac{\partial \widetilde{u}_1}{\partial y}=\dfrac{\lambda^3+2\lambda c^2\alpha^2}{2\,c\,(\lambda^2+c^2\alpha^2)}\;\widetilde{u}_1-\dfrac{\lambda^3}{2\,c\,(\lambda^2+c^2\alpha^2)}\;\widetilde{u}_3\\[0.85em]
\dfrac{\partial \widetilde{u}_3}{\partial y}=\dfrac{\lambda^3}{2\,c\,(\lambda^2+c^2\alpha^2)}\;\widetilde{u}_1-\dfrac{\lambda^3+2\lambda c^2\alpha^2}{2\,c\,(\lambda^2+c^2\alpha^2)}\;\widetilde{u}_3.
\end{cases}
\end{equation}

\begin{equation}\label{stripboundarycon}\widetilde{u}_1(\alpha,L\mathord{;}z)=0,\qquad \widetilde{u}_3(\alpha,0\mathord{;}z)=e^{i\alpha z}.\end{equation}

\begin{thm}\label{tildeuthm}
For $j=1,3$, the Fourier transforms $\widetilde{u}_j(\alpha,y\mathord{;}z)$ read
\begin{equation}\label{uhat1solu}\widetilde{u}_1(\alpha,y\mathord{;}z)=\frac{2\lambda^2\sinh\left(\frac{\lvert\alpha\lvert\lambda(L-y)}{\sqrt{\lambda^2+\alpha^2c^2}}\right)e^{i\alpha z}}{\big(\sqrt{\lambda^2+\alpha^2c^2}+\lvert\alpha\lvert c\big)^2\;e^{\frac{\lvert\alpha\lvert\,\lambda\,L}{\sqrt{\lambda^2+\alpha^2c^2}}}-\big(\sqrt{\lambda^2+\alpha^2c^2}-\lvert\alpha\lvert c\big)^2\;e^{-\frac{\lvert\alpha\lvert\,\lambda\,L}{\sqrt{\lambda^2+\alpha^2c^2}}}}\end{equation}
and
\begin{align}\widetilde{u}_3(\alpha,&\,y\mathord{;}z)=\nonumber\\
&e^{i\alpha z}\,\frac{\big(\sqrt{\lambda^2+\alpha^2c^2}+\lvert\alpha\lvert c\big)^2\;e^{\frac{\lvert\alpha\lvert\,\lambda\,(L-y)}{\sqrt{\lambda^2+\alpha^2c^2}}}-\big(\sqrt{\lambda^2+\alpha^2c^2}-\lvert\alpha\lvert c\big)^2\;e^{-\frac{\lvert\alpha\lvert\,\lambda\,(L-y)}{\sqrt{\lambda^2+\alpha^2c^2}}}}{\big(\sqrt{\lambda^2+\alpha^2c^2}+\lvert\alpha\lvert c\big)^2\;e^{\frac{\lvert\alpha\lvert\,\lambda\,L}{\sqrt{\lambda^2+\alpha^2c^2}}}-\big(\sqrt{\lambda^2+\alpha^2c^2}-\lvert\alpha\lvert c\big)^2\;e^{-\frac{\lvert\alpha\lvert\,\lambda\,L}{\sqrt{\lambda^2+\alpha^2c^2}}}}.\label{uhat3solu}\end{align}
For $j=0,2$, the Fourier transforms $\widetilde{u}_0(\alpha,y\mathord{;}z)$ and $\widetilde{u}_2(\alpha,y\mathord{;}z)$ follow immediately by the identities in equation (\ref{algebrain}).\end{thm}

\begin{proof}
We start by expressing the system (\ref{bivFouriersys2}) in matrix form. With the notation
$$A=\begin{pmatrix}
\dfrac{\lambda^3+2\lambda c^2\alpha^2}{2\,c\,(\lambda^2+c^2\alpha^2)}\; & \; -\dfrac{\lambda^3}{2\,c\,(\lambda^2+c^2\alpha^2)}\\
\dfrac{\lambda^3}{2\,c\,(\lambda^2+c^2\alpha^2)}\; & \; -\dfrac{\lambda^3+2\lambda c^2\alpha^2}{2\,c\,(\lambda^2+c^2\alpha^2)}
\end{pmatrix}$$
we can write that
\begin{equation}\label{smallvi}\frac{\partial}{\partial y}\begin{pmatrix}\widetilde{u}_1(\alpha,y\mathord{;}z)\\[0.5em] \widetilde{u}_3(\alpha,y\mathord{;}z)\end{pmatrix}=A\cdot\begin{pmatrix}\widetilde{u}_1(\alpha,y\mathord{;}z)\\[0.5em] \widetilde{u}_3(\alpha,y\mathord{;}z)\end{pmatrix}.\end{equation}
By using standard theory of linear systems of ordinary differential equations, the general solution to the system (\ref{smallvi}) can be represented in terms of the eigenvalues and eigenvectors of the matrix $A$. In particular, denote by $\theta_j,\;j=0,1,$ the eigenvalues of $A$, and let $\mathbf{v}_j,\;j=0,1,$ be the corresponding eigenvectors. We can write that 
$$\begin{pmatrix}\widetilde{u}_1(\alpha,y\mathord{;}z)\\[0.5em] \widetilde{u}_3(\alpha,y\mathord{;}z)\end{pmatrix}=\sum_{j=0}^1k_j\,e^{\theta_j y}\mathbf{v}_j$$ where $k_j,\;j=0,1$, are arbitrary constants. Straightforward calculations permit us to prove that
$$\theta_0=\frac{\lvert\alpha\lvert\,\lambda}{\sqrt{\lambda^2+\alpha^2c^2}},\qquad \theta_1=-\frac{\lvert\alpha\lvert\,\lambda}{\sqrt{\lambda^2+\alpha^2c^2}}.$$ The corresponding eigenvectors are
$$\mathbf{v}_0=\begin{pmatrix}\lvert\alpha\lvert c+\sqrt{\lambda^2+\alpha^2c^2}\\[0.75em]-\lvert\alpha\lvert c+\sqrt{\lambda^2+\alpha^2c^2}\end{pmatrix},\qquad \mathbf{v}_1=\begin{pmatrix}\lvert\alpha\lvert c-\sqrt{\lambda^2+\alpha^2c^2}\\[0.75em]-\lvert\alpha\lvert c-\sqrt{\lambda^2+\alpha^2c^2}\end{pmatrix}.$$
Thus, the general solution to the system (\ref{bivFouriersys2}) reads
$$\widetilde{u}_1(\alpha,y\mathord{;}z)=k_0\,(\sqrt{\lambda^2+\alpha^2c^2}+\lvert\alpha\lvert c)\;e^{\frac{\lvert\alpha\lvert\,\lambda\,y}{\sqrt{\lambda^2+\alpha^2c^2}}}-k_1\,(\sqrt{\lambda^2+\alpha^2c^2}-\lvert\alpha\lvert c)\;e^{-\frac{\lvert\alpha\lvert\,\lambda\,y}{\sqrt{\lambda^2+\alpha^2c^2}}}$$
and
$$\widetilde{u}_3(\alpha,y\mathord{;}z)=k_0\,(\sqrt{\lambda^2+\alpha^2c^2}-\lvert\alpha\lvert c)\;e^{\frac{\lvert\alpha\lvert\,\lambda\,y}{\sqrt{\lambda^2+\alpha^2c^2}}}-k_1\,(\sqrt{\lambda^2+\alpha^2c^2}+\lvert\alpha\lvert c)\;e^{-\frac{\lvert\alpha\lvert\,\lambda\,y}{\sqrt{\lambda^2+\alpha^2c^2}}}$$
By using now the boundary conditions (\ref{stripboundarycon}), it can be verified that
$$k_0=\frac{\big(\sqrt{\lambda^2+\alpha^2c^2}-\lvert\alpha\lvert c\big)\;e^{i\alpha z-\frac{\lvert\alpha\lvert\lambda L}{\sqrt{\lambda^2+\alpha^2c^2}}}}{\big(\sqrt{\lambda^2+\alpha^2c^2}-\lvert\alpha\lvert c\big)^2\;e^{-\frac{\lvert\alpha\lvert\,\lambda\,L}{\sqrt{\lambda^2+\alpha^2c^2}}}-\big(\sqrt{\lambda^2+\alpha^2c^2}+\lvert\alpha\lvert c\big)^2\;e^{\frac{\lvert\alpha\lvert\,\lambda\,L}{\sqrt{\lambda^2+\alpha^2c^2}}}}$$
and
$$k_1=\frac{\big(\sqrt{\lambda^2+\alpha^2c^2}+\lvert\alpha\lvert c\big)\;e^{i\alpha z+\frac{\lvert\alpha\lvert\lambda L}{\sqrt{\lambda^2+\alpha^2c^2}}}}{\big(\sqrt{\lambda^2+\alpha^2c^2}-\lvert\alpha\lvert c\big)^2\;e^{-\frac{\lvert\alpha\lvert\,\lambda\,L}{\sqrt{\lambda^2+\alpha^2c^2}}}-\big(\sqrt{\lambda^2+\alpha^2c^2}+\lvert\alpha\lvert c\big)^2\;e^{\frac{\lvert\alpha\lvert\,\lambda\,L}{\sqrt{\lambda^2+\alpha^2c^2}}}}.$$
Formulas (\ref{uhat1solu}) and (\ref{uhat3solu}) follow immediately.
\end{proof}

\end{document}
